\chapter{La herencia}
\label{ch:inheritance}

\section{Un país que no era pobre, y no estaba desarrollado}

El último día de 1958 Cuba tenía unos 6,6 millones de habitantes ---una quinta
parte de la población de España en una isla algo mayor que Portugal--- y era,
según casi todos los indicadores agregados, uno de los cuatro o cinco países más
ricos de América Latina.\unv{} Tenía más televisores por habitante que cualquier
país de la región y de los índices más altos del mundo; tenía una de las redes
ferroviarias más densas del hemisferio, heredada del azúcar; tenía
aproximadamente un médico por cada mil habitantes, una proporción comparable a
la de buena parte de Europa occidental de entonces; tenía una mayoría
alfabetizada, el 76 por ciento de la población mayor de diez años según el censo
de 1953; y tenía una esperanza de vida al nacer de alrededor de 62 años, superior
a la de todos los países latinoamericanos continentales salvo Argentina y
Uruguay \citep{perezjr2015}.\unv{}

Cada una de esas cifras es un promedio nacional, y cada una oculta lo que de
verdad importaba: que Cuba en 1958 no era un país sino dos economías atornilladas
una a la otra en el puerto de La Habana.

La economía urbana, y sobre todo la capital, era genuinamente moderna. La Habana
tenía grandes almacenes, un sistema de ómnibus con aire acondicionado, una bolsa
de valores, médicos especialistas, una prensa profesional de una docena de
diarios, una universidad con tradición investigadora real, y una clase media de
profesionales, comerciantes y obreros industriales sindicalizados cuyo nivel de
vida no habría desentonado en el sur de Europa. El movimiento obrero cubano era
fuerte ---los sindicatos azucarero y del transporte en particular--- y había
arrancado protecciones de empleo y salarios mínimos inusuales en la región, que
una década más tarde serían una de las restricciones con las que la revolución
tuvo que negociar en vez de rodear.

La economía rural era otra cosa. La fuente más citada sobre ella es una encuesta
a mil familias de trabajadores agrícolas realizada en 1956--57 por la Agrupación
Católica Universitaria, una organización laica sin simpatía alguna por lo que
vino después, que es la razón de que sus cifras hayan resistido el escrutinio.
Encontró que entre los trabajadores rurales el 4 por ciento comía carne con
alguna regularidad, el 1 por ciento comía pescado y el 11 por ciento tomaba
leche; que el 43 por ciento no sabía leer; que alrededor del 60 por ciento vivía
en viviendas de piso de tierra y techo de guano; que menos del 10 por ciento
tenía electricidad y cerca del 2 por ciento agua corriente
\citep{perezstable2012}; y que el ingreso anual de la familia mediana era una
fracción del salario urbano. La encuesta se publicó en una revista católica en
1957 y sus hallazgos no fueron seriamente rebatidos entonces. Aproximadamente el
43 por ciento de los cubanos vivía en el campo.

La descripción honesta de la Cuba prerrevolucionaria no es entonces ni la colonia
saqueada de las historias oficiales ni la república próspera de las memorias del
exilio. Es un país de ingreso medio con una capital de primer mundo, un interior
de tercer mundo, y una distribución tan ancha que el promedio nacional no
describe a casi nadie.

\section{El azúcar, y la forma que le dio a todo}

La razón de que existieran las dos economías es el azúcar, y la razón de que el
azúcar dominara es un instrumento de política en Washington llamado la cuota.

El azúcar era alrededor de cuatro quintas partes de los ingresos por
exportación cubanos, y lo había sido durante medio siglo.\unv{} Bajo los
acuerdos de comercio recíproco que regían el intercambio cubano-americano,
Estados Unidos compraba un tonelaje garantizado a un precio por encima del
mercado mundial y, a cambio, los aranceles cubanos favorecían las manufacturas
norteamericanas. El arreglo era, a corto plazo, extraordinariamente bueno para
Cuba: convertía una materia prima volátil en una anualidad estable y sostenía el
consumo de la clase media urbana. A largo plazo hizo tres cosas que se
acumularon.

Suprimió la industria. ¿Para qué levantar una fábrica cubana de zapatos si el
arancel abarataba los zapatos americanos y la cuota volvía ciertos los dólares
del azúcar? La manufactura cubana fuera del procesamiento azucarero, el tabaco y
un puñado de bienes de consumo era delgada, y lo que existía era en buena medida
ensamblaje.

Congeló la tierra. El azúcar quiere fincas grandes y planas y mano de obra
estacional barata, y obtuvo las dos. Hacia el censo agrícola de 1946, algo así
como el 8 por ciento de las fincas concentraba cerca del 70 por ciento de la
tierra agrícola, y una parte sustancial de la tierra cañera ---las estimaciones
habituales van de una cuarta parte a dos quintos de la capacidad de molienda---
pertenecía a compañías estadounidenses.\unv{} Buena parte de la mejor tierra se
mantenía enteramente fuera de producción como reserva para un buen año, en un
país que importaba un tercio de su comida.

Y volvió estructural el desempleo. La zafra corría de enero a mayo,
aproximadamente. El resto del año era el tiempo muerto, en el que varios cientos
de miles de cortadores de caña no tenían trabajo alguno. El desempleo medio anual
a mediados de los cincuenta rondaba el 16 por ciento y la oscilación estacional
era mucho mayor.\unv{} Un país no puede construir una política estable sobre una
economía que despide a la quinta parte de su fuerza laboral cada junio.

A esto se sumaba la cuestión de la propiedad, que importaba menos en lo económico
que en lo político. El capital estadounidense poseía la mayoría de las empresas
eléctricas y telefónicas, casi todo el níquel, una parte grande de la banca y las
refinerías, y la capacidad de molienda ya señalada. Si eso era extracción o
inversión es una pregunta genuinamente discutible ---las eléctricas eran
eléctricas de verdad y los centrales eran centrales de verdad---. Lo que no es
discutible es que ponía las alturas dominantes de la economía cubana a merced de
decisiones tomadas en otro país, y que eso resultaba intolerable para una
generación de cubanos criada en el nacionalismo de la revolución de 1933 y en la
constitución de 1940.

\section{La república que no pudo sostenerse}

La constitución cubana de 1940 era, sobre el papel, uno de los documentos más
avanzados de su época en cualquier parte: garantizaba la jornada de ocho horas,
las vacaciones pagadas, el derecho de huelga, igual salario para las mujeres y la
educación gratuita y obligatoria, y prohibía el latifundio y ponía límites a la
propiedad extranjera de la tierra. Muy poco del programa social llegó a
implementarse mediante leyes complementarias, y la cláusula sobre el latifundio
se ignoró sin más. Esa brecha entre una constitución excelente y un gobierno que
no la ejecutaba es el hecho central de la política republicana cubana, y es la
razón de que la constitución de 1940 se convirtiera en la bandera de toda
oposición posterior, la de Fidel Castro incluida.

Los gobiernos de Ramón Grau San Martín (1944--48) y Carlos Prío Socarrás
(1948--52) fueron electos, no fueron dictaduras, y fueron espectacularmente
corruptos. El dinero público desaparecía a una escala que se discutía
abiertamente en la prensa y que jamás se persiguió. La violencia política entre
facciones armadas, los grupos de acción, era rutinaria en La Habana. El Partido
Ortodoxo de Eduardo Chibás construyó un seguimiento de masas sobre una plataforma
de un solo tema, la honradez pública ---su lema radial era \emph{vergüenza contra
dinero}---, y cuando Chibás se disparó en pleno programa en agosto de 1951 el país
perdió a la única figura que quizá habría podido canalizar la corriente
anticorrupción hacia un resultado electoral. Castro era candidato ortodoxo a
representante en las elecciones que nunca se celebraron.

El 10 de marzo de 1952, tres meses antes de esas elecciones, Fulgencio Batista
---que había gobernado Cuba desde atrás y luego desde delante entre 1933 y 1944,
y que iba tercero en las encuestas--- tomó el poder en un golpe incruento al
amanecer. Esa es la bisagra. El golpe destruyó la vía constitucional, y la
destruyó contra un presidente que era corrupto pero electo y contra una
constitución de la que el país estaba orgulloso. Todo argumento insurreccional
formulado en Cuba entre 1952 y 1959 arranca ahí, y el argumento era correcto:
después de marzo de 1952 no había forma legal de remover al gobierno.

El segundo período de Batista fue una dictadura, y empeoró con el tiempo. El
Congreso fue suspendido y luego restaurado como mero sello ratificador mediante unas
elecciones amañadas en 1954 que la oposición boicoteó; la prensa fue censurada
por oleadas; y, una vez comenzada la insurgencia guerrillera, la policía y el
ejército respondieron con tortura, desaparición y asesinatos públicos, sobre todo
por parte de la policía habanera bajo Esteban Ventura y del ejército en Oriente.
El número de personas que el régimen mató entre 1952 y 1958 es genuinamente
desconocido. La cifra de 20.000, que entró en circulación a través de una revista
habanera en enero de 1959 y que luego se repitió durante sesenta años en los
textos oficiales cubanos, nunca ha sido respaldada por ningún conteo y no la cree
ningún historiador serio de ningún bando; las estimaciones cuidadosas van de unos
dos mil a unos cinco mil, combatientes incluidos \citep{thomas1971}.\unv{} Es una
cifra menor y enteramente suficiente.

Entretanto La Habana se convirtió en el distrito de placer extraterritorial de
Estados Unidos. Casinos había tenido antes; bajo Batista, y al amparo de una ley
explícita de 1955 que ofrecía licencias de juego y exenciones fiscales a quien
construyera hoteles, el crimen organizado norteamericano ---Meyer Lansky y Santo
Trafficante en primer lugar--- pasó a ser parte estructural de la economía
turística, con el régimen cobrando su parte. Aquí es donde tiene que empezar el
capítulo sobre turismo de la Parte III, porque la memoria de aquella Habana sigue
haciendo trabajo político en Cuba setenta años después, y toda propuesta de
convertir el turismo en la industria ancla del país tiene que responder a ella.

\section{La raza, que la república no resolvió}

El censo de 1953 registró que alrededor del 27 por ciento de los cubanos se
identificaba como negro o mestizo, cifra que la mayoría de los demógrafos
considera un subregistro.\unv{} La esclavitud había terminado en 1886, más tarde
que en cualquier lugar de las Américas salvo Brasil, y dentro de la memoria viva
de personas que estaban vivas en 1958. La república había abolido la distinción
racial legal y había producido un ejército, una policía y una clase política
genuinamente multirraciales ---el propio Batista era mestizo y se le vetaba la
entrada a los clubes sociales blancos de La Habana mientras gobernaba el país---.

Lo que no había abolido era la barrera de color privada: las playas, los clubes,
los hoteles y los barrios que eran blancos en la práctica; la contratación en la
banca, el comercio y las profesiones que era blanca por preferencia; la forma de
la escalera laboral, en la que los afrocubanos se concentraban en el corte de
caña, la construcción, el servicio doméstico y la economía informal. Un intento
de legislar contra la discriminación en el empleo a finales de los cuarenta
fracasó \citep{delafuente2001}. Cuando la revolución desegregó las playas y los
clubes en 1959 estaba atendiendo algo real, lo hizo rápido, y se ganó entre los
afrocubanos una lealtad que ha sobrevivido a casi cualquier otra fuente de apoyo
del gobierno ---hecho que la Parte III tiene que tomarse en serio, porque el
sector social más expuesto a una transición mal diseñada es el mismo---.

\section{La insurrección}

La oposición armada a Batista no fue una sola organización, y la de Castro no fue
inicialmente la mayor.

Empezó el 26 de julio de 1953 con el ataque de unos 130 jóvenes ortodoxos al
cuartel Moncada en Santiago, que fracasó por completo; cerca de la mitad de los
atacantes murió, la mayoría después de ser capturada. La defensa de Castro en el
juicio, publicada más tarde como \emph{La historia me absolverá}, expuso un
programa ---restaurar la constitución de 1940, reforma agraria, reparto de
utilidades, educación, nacionalizar los servicios públicos--- que era
nacionalista radical y no comunista, y que los cubanos podían leer, y leyeron,
como una continuación de Chibás. Fue condenado, amnistiado en 1955 en un gesto de
confianza de un régimen que no esperaba lo que vino después, y se marchó a
México.

El \emph{Granma} desembarcó 82 hombres en Oriente el 2 de diciembre de 1956; una
docena llegó a la Sierra Maestra. Desde allí el Movimiento 26 de Julio libró una
pequeña guerra de guerrillas ---nunca más de unos pocos cientos de combatientes
en las montañas--- mientras una clandestinidad urbana mucho mayor, el llano,
manejaba el sabotaje, las finanzas, la propaganda y el abastecimiento en las
ciudades bajo Frank País, hasta que la policía lo mató en julio de 1957
\citep{sweig2002}. Una organización aparte, el Directorio Revolucionario de los
estudiantes universitarios, estuvo a punto de matar a Batista en el asalto al
Palacio Presidencial de marzo de 1957 y perdió en ello a casi toda su dirigencia.
El viejo partido comunista, el Partido Socialista Popular, se opuso a la lucha
armada hasta tarde y solo se sumó en 1958, hecho que importaría enormemente en
los años sesenta.

Dos cosas decidieron la guerra, y ninguna fue militar en el sentido corriente. La
primera es que la huelga general convocada por el movimiento en abril de 1958
fracasó, lo que desplazó el centro de gravedad de manera decisiva de las ciudades
a la Sierra y de la dirección del llano a la de Castro. La segunda es que en
marzo de 1958 Estados Unidos, presionado por el Congreso a causa de las
atrocidades del régimen, embargó la venta de armas a Batista. El ejército no fue
tanto derrotado en el campo como que dejó de ser un ejército: la ofensiva del
verano de 1958 sobre la Sierra se derrumbó, unidades enteras se rindieron o
cambiaron de bando, y cuando la columna del Che Guevara tomó Santa Clara a fines
de diciembre el camino a La Habana quedó abierto. Batista voló a la República
Dominicana en la madrugada del 1 de enero de 1959.

\section{Lo que de verdad se heredó}

Conviene enunciarlo con claridad, porque las historias posteriores lo distorsionan
en ambas direcciones.

\begin{ledger}{la herencia, 1 de enero de 1959}
\emph{Activos.} Una economía de ingreso medio con infraestructura funcional, una
mayoría urbana alfabetizada, una clase profesional real, sindicatos fuertes, una
red vial y ferroviaria densa, una industria azucarera sustancial con centrales
modernos, recursos minerales y una planta turística que ya funcionaba. Una
constitución con legitimidad de masas que todos decían estar restaurando. Deuda
externa baja.

\emph{Pasivos.} Una economía exportadora de monocultivo dependiente de la
voluntad política de un solo comprador; desempleo estacional estructural cercano
a la quinta parte de la fuerza laboral; una mitad rural del país en niveles
haitianos de nutrición, alfabetización y vivienda; propiedad de la tierra
concentrada y buena parte de la mejor tierra ociosa; propiedad extranjera de los
servicios públicos, las minas y gran parte de la molienda; una barrera de color
sin resolver; una cultura política en la que la corrupción se daba por supuesta y
la vía constitucional acababa de cerrarse por la fuerza; y ninguna institución
con legitimidad para arbitrar lo que viniera después ---el ejército era de
Batista, el congreso era una ficción, los tribunales estaban comprometidos y los
partidos eran irrelevantes desde 1952---.
\end{ledger}

Ese último pasivo es el que decidió la forma de todo lo que siguió. En enero de
1959 el único cuerpo en Cuba con legitimidad inequívoca era el movimiento
guerrillero y, dentro de él, un hombre. No había contrapeso, y no lo había desde
el 10 de marzo de 1952. Mucho de lo que sigue en este libro ---la velocidad de la
radicalización, la ausencia de todo freno interno, la facilidad con que una
coalición revolucionaria plural se convirtió en un partido único--- es
consecuencia de que el golpe de Batista ya había destruido las instituciones que
podrían haberla frenado. La revolución no se apoderó de un orden constitucional.
Heredó un vacío donde había habido uno.

Esta es la primera lección que el libro lleva a la Parte III, y no es abstracta:
\emph{la secuencia en que se reconstruyen las instituciones importa más que la
velocidad}, y una transición que produzca un actor legítimo y ningún otro
terminará donde terminó 1959, cualesquiera que sean sus intenciones. Cuba ya ha
corrido ese experimento dos veces, en 1952 y en 1959, desde direcciones opuestas.
